\histitem{Connection between harmonic analysis and group theory}

We have described, in the historical notes of the Book on Integration (cf. ÉHM, p. 275), how questions related to the heat equation led Fourier to the revolutionary conception of representing an arbitrary function as a sum of trigonometric functions.

We cannot trace here the development of the theory of Fourier series and integrals throughout the nineteenth century, nor the profound influence that the questions raised by this theory had on the evolution of Analysis --- to the point of overturning the conception of real numbers, contributing to the birth of set theory (cf. ÉHM, p. 42). But in order to understand the form of the ideas set forth in this Book, it is appropriate to describe the manner in which harmonic analysis was connected after 1925 to group theory and to Hilbertian spectral theory.

Recognition of the link between Fourier's ideas and the group structures of R/Z and R, as well as the generalization of these ideas to other groups, came rather late.

One can naturally, in the early stages of the theory of finite groups and their characters, identify results that today appear to contain the rudiments of finite Fourier analysis. For example, the dual of the group of invertible elements of Z/qZ, as well as the orthogonality relation expressing the characteristic function of an element as a linear combination of characters, appear implicitly in 1837 in Dirichlet's article {[}17{]} proving that there are infinitely many prime numbers \(p \equiv a \mod q\) if a is relatively prime to q. He relies on ideas of Gauss, who had introduced the term “character” in his study of binary quadratic forms with integer coefficients and fixed discriminant (cf. {[}22, § 230{]}).

But neither Gauss nor Dirichlet uses the language of groups. It was Dedekind, presenting these works in 1879, who recognized their latent role and defined the notion of character for commutative groups. Dedekind's remarks were considerably amplified by H. Weber, first in 1882 {[}77{]}, then in 1886 {[}78, p. 112{]}, where he introduced the dual group of a finite abelian group A and noted that the latter is isomorphic to A. His purpose was then to construct abelian extensions of the field Q with given Galois group, and he did not consider the problem of harmonic analysis on such a group.

In the case of noncommutative groups, the foundations of the theory of representations of finite groups were solidly established around the turn of the twentieth century, following the works of G. Frobenius, W. Burnside, and I. Schur (cf. ÉHM, p. 154). After 1905, it was clear that the orthogonality relations of characters played a crucial role in the organization of the theory. But their kinship with harmonic analysis remained hidden.

It is to Hermann Weyl that we owe the connection between these fields, when he conceived the idea of a general theory of representations for compact groups. In several major texts published between 1925 and 1927, he recognized the links that would unite such a theory with Fourier analysis and Hilbertian spectral theory, and he laid the foundations for many later discoveries.

It was a letter from Schur that set Weyl on this path, thanks to their common interest in invariant theory. We described, in the note on Haar measure (ÉHM, pp. 289–291), how A. Hurwitz had conceived as early as 1897 the idea of using an “invariant integration” to determine the polynomials on \(R^{n}\) invariants under the orthogonal group. It seems that Schur was not aware of Hurwitz's results before 1924, when he suddenly understood how invariant integration could be used to extend to the group \(\mathbf{O}(n)\) character theory and the orthogonality relations he had established for finite groups. Weyl immediately saw that Schur's methods, combined with the work of Élie Cartan, made it possible to construct the irreducible representations of compact semisimple Lie groups. The series of three memoirs in which Weyl blended the ideas of Cartan and Schur, published in 1925 {[}83{]}, already contained the essential results of LIE, IX, §6–7; cf. ÉHM, p. 328–330.

But it was the following year, in an equally famous text written with his student F. Peter {[}54{]}, that Weyl brought to light the connections between the representation theory of compact groups, Fourier analysis, and Hilbertian spectral theory. This text of fewer than twenty pages contains in embryo a great many future developments.

Peter and Weyl free themselves of any algebraic hypothesis on the structure of the group under study, assuming only that G is a compact topological group equipped with an invariant measure. The existence of such a measure was clear for connected Lie groups; soon A. Haar would establish it in general, partly motivated by the results presented here (cf. ÉHM, p. 291).

The starting point of their study is the orthogonality of matrix coefficients. Schur had observed that by choosing a representative \(\pi \in \widehat{\mathrm{G}}\) for each equivalence class of irreducible representations, a G-invariant inner product on the space E of \(\pi\) , and an orthonormal basis of E, one obtains a family of matrix coefficients that is orthonormal in \(\mathrm{L}^2 (\mathrm{G})\) . Peter and Weyl intend to show that such a family is complete.

In the case of finite groups, the analogous result had been proved by Frobenius by studying the group algebra C{[}G{]}. Peter and Weyl then consider the space \(\mathcal{C}(\mathrm{G})\) of continuous functions on G and equip it with the convolution product. This appears to be the first use of the convolution product as an abstract operation in explicit connection with the group structure (cf. ÉHM, p. 295). Peter and Weyl moreover call Gruppenzahl ("group number") an element of the algebra \(\mathcal{C}(\mathrm{G})\) , and denote by xy the (convolution) product of two Gruppenzahlen. They also equip \(\mathcal{C}(\mathrm{G})\) with the involution \(f \mapsto \widetilde{f}\) , where \(\widetilde{f}(g) = \overline{f(g^{-1})}\) for \(g \in G\) .

Having thus introduced the involutive algebra \(\mathcal{C}(\mathrm{G})\) , the first observation made by Peter and Weyl is that every unitary representation \(\pi\) of G gives rise to a representation \(f \mapsto \pi(f)\) of \(\mathcal{C}(\mathrm{G})\) (cf. V, p. 400). The notion of Hermitian element of \(\mathcal{C}(\mathrm{G})\) is explicitly defined, as is (implicitly) that of positive element, which opens the way to the application of Hilbert space techniques.

Peter and Weyl no longer hesitate to call the operator \(\pi(f)\) the "Fourier coefficient" of \(f\) , and continue the parallel with Fourier theory by noting that Schur's orthogonality relations imply Bessel's inequality

\[
\sum_ {\pi \in \widehat {\mathrm{G}}} \dim (\pi) \operatorname{Tr} (\pi (f) \pi (f) ^ {*}) \leqslant (f * \widetilde {f}) (e) = \| f \| _ {2} ^ {2}.\tag{6}
\]

The Hilbert--Schmidt spectral theory then allows them to show that this is an equality ("Plancherel formula"). To every element \(\varphi\) of \(\mathcal{C}(\mathrm{G})\) , Peter and Weyl associate the continuous kernel \(K_{\varphi}\colon G\times G\to C\) defined by \(\mathrm{K}_{\varphi}(x,y)=\varphi(xy^{-1})\) . They then observe that if \(\varphi\) is of the form \(f*\widetilde{f}\) with \(f\in\mathcal{C}(\mathrm{G})\) , \(^{(6)}\) then \(K_{\varphi}\) is a positive Hermitian kernel in the sense of Hilbert--Schmidt theory. A variant of a Schmidt algorithm \(^{(7)}\) then allows them to construct an orthonormal basis of eigenfunctions of the kernel \(K_{\varphi}\) , and then to reduce the proof of the equality in (6) to the fact that the trace of the operator defined by \(K_{\varphi}\) is equal to the sum of its eigenvalues.

Naturally, only the representations \(\pi\) satisfying \(\pi(f) \neq 0\) can occur in (6); the spectral theory for \(K_{\varphi}\) shows the existence of such a representation if \(f\) is not identically zero. Peter and Weyl easily deduce from this that the irreducible representations separate the points of G --- a special case of the Gelfand--Raikov theorem, which will be discussed later.

In the celebrated works of Frobenius, an equality analogous to (6) made it possible to show that the regular representation of G contains all irreducible representations. But this involved applying this equality by taking for f the identity element of C{[}G{]}; one then obviously had \(\pi(f) \neq 0\) for all \(\pi\) . According to Peter and Weyl, the absence of an identity element for convolution explains many difficulties in their proof; they nevertheless observe, borrowing from the theory of Fourier series an idea destined for a great future, that one can deduce the analogue of Frobenius's result by applying (6) to a sequence of functions forming an approximate identity for convolution (cf. no. 10 of I, p. 120).

Likewise, starting from the equality arising from (6), the polarized version

\[
\sum_ {\pi \in \widehat {\mathrm{G}}} \dim (\pi) \operatorname{Tr} (\pi (x) \pi (y)) = (x * y) (e)\tag{7}
\]

makes it possible to obtain, by taking for \(y\) an approximate identity for convolution, a uniform approximation of every element of \(\mathcal{C}(\mathrm{G})\) by linear combinations of matrix coefficients of irreducible representations. These methods had just enabled Weyl to give a new proof {[}84{]} of the fundamental result of the theory of "almost periodic functions" of H. Bohr, to which we shall return shortly. Peter and Weyl point out that one can read therein the first analytic application of the theory of representations of a noncompact group, here that of the translations of \(\mathbf{R}\) .

The young quantum theory very quickly provided Weyl with the opportunity to return to the representations of R. At the end of 1927, he noted the interest of groups (finite or not) and of representations for clarifying the foundations of this theory {[}85{]}. He would return to this the following year, in a resounding book {[}86{]}.

In his 1927 article, Weyl observes that if u is a continuous self-adjoint operator on a Hilbert space, then \(t \mapsto e^{itu}\) is a unitary representation of R. But there exist unitary representations of R that are not of this form; taking up von Neumann's idea on the role of symmetric partial operators in representing observable physical quantities, Weyl suggests that every unitary representation of R is of the form \(t \mapsto e^{itu}\) where u is a partial self-adjoint operator on a Hilbert space. This point of view allows him to reduce the analysis of the "canonical relations" of Heisenberg and Schrödinger, already studied in detail by von Neumann, to that of the connections between two unitary representations of R on the same Hilbert space. Stone {[}73{]} shows in 1932 the result hoped for by Weyl (V, p. 428, th. 1), in the wake of his study of the spectral properties of self-adjoint operators.

